% $Id: header.tex 1229 2009-10-23 13:58:42Z inf6254 $
%%%%%%%%%%%%%%%%%%%%%%%%%%%%%%%%%%%%%%%%%%%%%%%%%%%%%%%%%%%%%%%%%%%%%%%
\documentclass
  [ %twoside            % beidseitiger Druck
  , openright           % Kapitel beginnen auf einer rechten Seite
  , listof=totoc        % Verzeichnisse im Inhaltsverzeichnis
  , bibliography=totoc  % Literaturverzeichnis im Inhaltsverzeichnis
  , parskip=half        % Absätze durch einen vergrößerten Zeilenabstand getrennt
%  , draft              % Entwurfsversion
  ]{scrreprt}           % Dokumentenklasse: KOMA-Script Buch

\bibliographystyle{alpha}


%%%%%%%%%%%%%%%%%%%%%%%%%%%%%%%%%%%%%%%%%%%%%%%%%%%%%%%%%%%%%%%%%%%%%%%
% Packages
%%%%%%%%%%%%%%%%%%%%%%%%%%%%%%%%%%%%%%%%%%%%%%%%%%%%%%%%%%%%%%%%%%%%%%%
\usepackage{scrhack}

 
\usepackage{ifpdf}
\ifpdf
  \usepackage{ae}               % Fonts für pdfLaTeX, falls keine cm-super-Fonts installiert
  \usepackage{microtype}        % optischer Randausgleich, falls pdflatex verwandt
  \usepackage[pdftex]{graphicx} % Grafiken in pdfLaTeX
\else
  \usepackage[dvips]{graphicx}  % Grafiken und normales LaTeX
\fi

\usepackage[a4paper,
            inner=2.0cm,outer=2.0cm,bindingoffset=0.5cm,
            top=1.5cm,bottom=1.5cm,
            footskip=1.0cm,includeheadfoot]{geometry}
            
\usepackage[utf8]{inputenc}         % Input encoding (allow direct use of special characters like "ä")
%\usepackage[english]{babel}
\usepackage[ngerman]{babel}
\usepackage[T1]{fontenc}
\usepackage[automark]{scrlayer-scrpage} % Schickerer Satzspiegel mit KOMA-Script
\usepackage{setspace}           	    % Allow the modification of the space between lines
\usepackage{booktabs}           	    % Netteres Tabellenlayout
\usepackage{multicol}                   % Mehrspaltige Bereiche
\usepackage{quotchap}                   % Beautiful chapter decoration
\usepackage{acronym}                    % list of acronyms and abbreviations
\usepackage{subfig}                     % allow sub figures
\usepackage{tabularx}                   % Tabellen mit fester Breite
\usepackage{lipsum}                     % für Blindtexte
\usepackage{titlesec}                   % für eigene Überschriftenstile
\usepackage{nameref}                    % Namen von refs bekommen (Kapitel etc.)
\usepackage{todonotes}                  % ToDo-Kommentare und Auflistung
\usepackage{float}                      % Bilder positionieren
\usepackage{interval}                   % Intervalle
\usepackage{amssymb}                    % Mathekrams, z.B. \mathbb für Mengen
\usepackage{amsmath}                    % Noch mehr Mathekrams, für argmax etc.
\usepackage[ruled,vlined]{algorithm2e}  % Pseudocode, Algorithmenbeschreibung
\usepackage{listings}                   % Nice listings with frames etc.
\usepackage{lstautogobble}              % Auto gobbling for listings.
\usepackage{stackengine}                % Stacking of objects
 


%%%%%%%%%%%%%%%%%%%%%%%%%%%%%%%%%%%%%%%%%%%%%%%%%%%%%%%%%%%%%%%%%%%%%%%
% User Commands
%%%%%%%%%%%%%%%%%%%%%%%%%%%%%%%%%%%%%%%%%%%%%%%%%%%%%%%%%%%%%%%%%%%%%%%

\newcommand*{\Name}[1]{\textsc{#1}}
\newcommand*{\Befehl}[1]{\texttt{#1}}
\newcommand*{\Fachbegriff}[1]{\emph{#1}}
\newcommand*{\FachbegriffMitAbk}[2]{\emph{#1} (#2)}
\newcommand*{\ErsterFachbegriffMitAbk}[3]{\emph{#1} (#2)\acused{#3}}
\newcommand*{\sogenannt}[1]{\glqq{}#1\grqq{}}
\newcommand*{\aufzaehlung}[1]{\glqq{}#1\grqq{}}
\newcommand{\hiddensection}[1]{ \stepcounter{section} \section*{#1}}
\newcommand*{\inlinecode}[1]{\texttt{#1}}
\newcommand*{\filepath}[1]{\texttt{#1}}


\DeclareMathOperator*{\mathmax}{max}
\DeclareMathOperator*{\mathmin}{min}
\DeclareMathOperator*{\argmax}{arg\,max}
\DeclareMathOperator*{\argmin}{arg\,min}


%%%%%%%%%%%%%%%%%%%%%%%%%%%%%%%%%%%%%%%%%%%%%%%%%%%%%%%%%%%%%%%%%%%%%%%
% Layout
%%%%%%%%%%%%%%%%%%%%%%%%%%%%%%%%%%%%%%%%%%%%%%%%%%%%%%%%%%%%%%%%%%%%%%%

\lstset{ % For custom listings
    numbers=left, 
    numberstyle=\small, 
    numbersep=15pt, 
    frame = single, 
    language=Java, 
    framexleftmargin=25pt,
    autogobble=true
}

\pagestyle{scrheadings}
% \pagestyle{empty}
\clubpenalty = 10000
\widowpenalty = 10000
\displaywidowpenalty = 10000

\makeatletter
\renewcommand{\fps@figure}{htbp}
\makeatother


%% Document properties %%%%%%%%%%%%%%%%%%%%%%%%%%%%%%%%%%%%%%%%%%%%%%%%
\newcommand{\projname}{Reinforcement-Utility-Village}

% Der Titel ist Banane, hier soll einfach irgendwas stehen
\newcommand{\titel}{Anwendung von Reinforcement Learning zur Entwicklung authentischer Spielwelten mithilfe bedürfnisbasierter Agenten}
\newcommand{\authorname}{Tobias Jansing}
\newcommand{\thesisname}{Master Thesis}
\newcommand{\untertitel}{Ein alternativer Ansatz zur Entwicklung von \textit{Non-Player-Characters} anhand einer Dorfsimulation mit Unity}
\newcommand{\Datum}{26 Februar 2020}

\ifpdf
  \usepackage{hyperref}
  \definecolor{darkblue}{rgb}{0,0,.5}
  \hypersetup
  	{ colorlinks=true
  	, breaklinks=true
    , linkcolor=darkblue
    , menucolor=darkblue
    , urlcolor=darkblue
    , citecolor=darkblue
    , pdftitle={\projname -- \untertitel}
    , pdfsubject={\thesisname}
    , pdfauthor={\authorname}
    }
\else
\fi

%% Listings %%%%%%%%%%%%%%%%%%%%%%%%%%%%%%%%%%%%%%%%%%%%%%%%%%%%%%%%%
\usepackage{listings}
\KOMAoptions{listof=totoc} % necessary because of scrhack
\renewcommand{\lstlistlistingname}{List of Listings}
\lstset
  { basicstyle=\small\ttfamily
  , breaklines=true
  , captionpos=b
  , showstringspaces=false
  , keywordstyle={}
  }

\lstnewenvironment{inlinehaskell}
{\spacing{1}\lstset{language=haskell,nolol,aboveskip=\bigskipamount}}
{\endspacing}

\lstnewenvironment{inlinexml}
{\spacing{1}\lstset{language=XML,nolol,aboveskip=\bigskipamount}}
{\endspacing}

\newcommand{\haskellinput}[2][]{
  \begin{spacing}{1}
  \lstinputlisting[language=Haskell,nolol,aboveskip=\bigskipamount,#1]{#2}
  \end{spacing}
}

\newcommand{\haskellcode}[2][]{\mylisting[#1,language=Haskell]{#2}}

\newcommand{\mylisting}[2][]{
\begin{spacing}{1}
\lstinputlisting[frame=lines,aboveskip=2\bigskipamount,#1]{#2}
\end{spacing}
}
